% ================================================================
% ==== FILE: content/k_levels/k12.tex
% ================================================================

% ==============================
%  Ontology of Continua — Core
%  K-level Module: K12 (Limit Continua)
%  FULL MODULE — FINAL
% ==============================

\subsubsection{$K_{12}$ Übersicht}
\label{sec:k12-overview}

$K_{12}$ ist das \emph{Grenzkontinuum}: der strukturelle Fixpunkt des
rekursive Metahierarchie, entwickelt in $K_{11}$.
Es ist die erste Ebene, in der:

\begin{itemize}
    \item rekursive Metastrukturen stabilisieren,
    \item-Verzweigung (Axiom~21) erreicht den Abschluss,
    \item ein globaler struktureller Fixpunkt existiert,
    \item alle Operatoren $(\Psi,\Phi,\Lambda,U,\Chi)$ wirken als Symmetrien,
    \item wird die gesamte Leiter $K_0\dots K_{11}$ innerhalb einer darstellbar
          invariantes Kontinuum.
\end{itemize}

$K_{12}$ ist daher nicht nur „die nächste Ebene“, sondern die strukturelle Grenze von
der K-Tower in Core v2.6.

% ================================================================
\subsubsection{Zustandsraum $\Omega(K_{12})$}

\[
\Omega(K_{12}) =
\{
\Omega^*,\;
A^*,\;
P^*,\;
\Theta^*,\;
J^*,\;
C^*,\;
\mu^{(12)}_t
\},
\]

wo:
\begin{itemize}
    \item $\Omega^*$ ist der Grenzzustandsraum, der alle zulässigen Einbettungen enthält
          von $K_0\!\to\!K_{11}$,
    \item $A^*$ ist die globale Achsenmenge, die durch die Schließung aller Achsen entsteht
          aus früheren Levels,
    \item $P^*$ sind Potentiale, die rekursive Konvergenz überleben,
    \item $\Theta^*$ sind stabile Schwellenwerte, die unter Rekursion erhalten bleiben,
    \item $J^*$ sind Flüsse, die unter allen Strukturoperatoren invariant bleiben,
    \item $C^*$ ist die feste Familie von Zyklen, die die Zeit von $K_{12}$ definiert,
    \item $\mu^{(12)}_t$ ist das Maß über Grenzkonfigurationen.
\end{itemize}

Alle Komponenten erfüllen die globale Invarianz:
\[
X^* = \mathcal{R}(X^*),\quad \forall X \in \{\Omega, A, P, \Theta, J, C\}.
\]

% ================================================================
\subsubsection{Grenze $\partial\Omega(K_{12})$}

Die Grenze besteht aus:
\begin{itemize}
    \item-Konfigurationen, bei denen die Rekursion nicht konvergiert,
    \item Punkte, an denen Verzweigungen zu inkompatiblen globalen Topologien führen,
    \item gibt an, dass die Grenzwertinvarianz unter Operatoren verletzt wird,
    \item Divergenz von Festkommazyklen,
    \item Unzulässigkeit in $M_{12}$.
\end{itemize}

Das Überschreiten von $\partial\Omega(K_{12})$ zerstört die Grenzstruktur und macht sie rückgängig
System zu einem scheiternden $K_{11}$-Regime.

% ================================================================
\subsubsection{Achsen $A(K_{12})$}

\[
A(K_{12}) =
\overline{\bigcup_{i=0}^{11} A(K_i)}
\quad \text{vorbehaltlich rekursiver Schließung}.
\]

Dieser Achsensatz ist:
\begin{itemize}
    \item geschlossen unter Rekursion von $K_{11}$,
    \item stabil unter allen Operatoren $(\Psi,\Phi,\Lambda,U,\Chi)$,
    \item minimal unter den Achsensätzen, die die globale Invarianz erfüllen.
\end{itemize}

Jede Achse wird zu einer \emph{Symmetrierichtung} in $K_{12}$.

% ================================================================
\subsubsection{Potentiale $P(K_{12})$}

\[
P(K_{12}) = \lim_{n\to\infty} P^{(n)}(K_{11}),
\]

d.h.\ die Potentiale, die bei Rekursion endlich, kohärent und geschlossen bleiben:
\begin{itemize}
    \item $P_{\mathrm{fix}}^*$ — Festkomma-Unterstützungspotential,
    \item $P_{\mathrm{coh}}^*$ — globales Kohärenzpotential,
    \item $P_{\mathrm{branch}}^*$ — stabilisiertes Verzweigungspotential,
    \item $P_{\mathrm{embed}}^*$ – universelles Einbettungspotential.
\end{itemize}

% ================================================================
\subsubsection{Schwellenwerte $\Theta(K_{12})$}

\[
\Theta(K_{12}) =
\lim_{n\to\infty} \Theta^{(n)}(K_{11}),
\]

geben:
\begin{itemize}
    \item $\Theta^*_{\mathrm{fix}}$ – Fixkomma-Existenzschwelle,
    \item $\Theta^*_{\mathrm{coh}}$ – globale Kohärenzschwelle,
    \item $\Theta^*_{\mathrm{branch}}$ – Verzweigungs-Schließungsschwelle,
    \item $\Theta^*_{\mathrm{inv}}$ — Invarianzschwelle für Operatoren,
    \item $\Theta^*_{\mathrm{collapse}}$ – strukturelle Kollapsgrenze.
\end{itemize}

Alle Schwellenwerte werden zu \emph{global stabilen Konstanten} für das Grenzkontinuum.

% ================================================================
\subsubsection{Flüsse $J(K_{12})$}

\[
J(K_{12}) = \{J^*\mid \Phi(J^*)=J^*\}.
\]

Somit muss jeder Fluss sein:
\begin{itemize}
    \item invariant unter rekursiver Evolution,
    \item kompatibel mit Grenzschwellen,
    \item global kohärent über Projektionen und Einbettungen hinweg,
    \item wurde unter Bedieneraktion behoben.
\end{itemize}

Dies ist die einzige Ebene, auf der \emph{alle Flüsse Symmetrieflüsse sind}.

% ================================================================
\subsubsection{Zyklen $C(K_{12})$}

\[
C(K_{12}) = \{ C^*\mid \Pi(C^*)>\Theta_{\mathrm{time}}^*,\;
\Lambda(C^*)=C^*\}.
\]

Dies sind die Zyklen, die:
\begin{itemize}
    \item bleibt bei allen rekursiven Verfeinerungen bestehen,
    \item definiert die Entstehungszeit von $K_{12}$,
    \item bleibt bei Verzweigungen und Bedieneraktionen unveränderlich,
    \item dienen als strukturelles Rückgrat des Kontinuums.
\end{itemize}

% ================================================================
\subsubsection{Effektive Zeit $\tau_{\mathrm{eff}}(K_{12})$}

Die Zeit wird durch invariante Zyklen bestimmt:
\[
\tau_{\mathrm{eff}}(K_{12}) = \max_j \tau_{C_j^*}.
\]

Eigenschaften:
\begin{itemize}
    \item alle zeitlichen Richtungen sind Symmetrien,
    \item keine neuen Zeitrichtungen entstehen können,
    \item-Zeit ist global kohärent und rekursiv unveränderlich.
\end{itemize}

% ================================================================
\subsubsection{Kontinuität $k(K_{12})$}

\[
k_{12} =
H(\Omega(K_{12}))
\cdot
\frac{|C^*|}{|C_{\max}|}
\cdot
\frac{|A^*|}{|A_{\max}|}
\cdot
\frac{P_{\mathrm{fix}}^*}{P_{\mathrm{fix,max}}}
\cdot
\left(1-\frac{T_{12}}{\Theta^*_{\mathrm{collapse}}}\right)_+.
\]

$k_{12}$ erreicht den maximal möglichen Wert unter allen $K$-Stufen:
\[
k_{12} = \max_{i=0\dots12} k(K_i).
\]

% ================================================================
\subsubsection{Strukturelle Spannung $T(K_{12})$}

Quellen:
\begin{itemize}
    \item ungelöste Verzweigungsinkompatibilitäten,
    \item-Operatorinvarianzfehler,
    \item Instabilität von Fixpunkten (selten, aber möglich),
    \item-Projektion – Einbettung von Inkonsistenzen,
    \item Verletzungen der $M_{12}$-Grenzen.
\end{itemize}

Ein Zusammenbruch tritt auf, wenn:
\[
T(K_{12}) > \Theta^*_{\mathrm{collapse}}.
\]

% ================================================================
\subsubsection{Energie $E(K_{12})$}

\[
E(K_{12}) =
E_{\mathrm{fix}}^*
+E_{\mathrm{coh}}^*
+E_{\mathrm{branch}}^*
+E_{\mathrm{embed}}^*
+E_{\mathrm{sym}}.
\]

Interpretation:
\begin{itemize}
    \item minimale Energie, die zur Erhaltung der Festkommastruktur erforderlich ist,
    \item Energie der globalen Kohärenzerhaltung,
    \item energiestabilisierender Verzweigungsverschluss,
    \item Energie, die universelle Einbettungen unterstützt,
    \item Symmetrieenergie der Operatorinvarianz.
\end{itemize}

$K_{12}$ ist das energiestabilste zulässige Kontinuum.

% ================================================================
\subsubsection{Operatoren auf $K_{12}$ ($\Psi$, $\Phi$, $\Lambda$, $U$, $\Chi$)}

Auf der Grenzebene:

\begin{itemize}
    \item $\Psi$ – wird Identität bei zulässigen Zuständen,
    \item $\Phi$ – wird zum globalen Invarianzfluss,
    \item $\Lambda$ – Schließung aller Grenzzyklen,
    \item $U$ – universeller Einbettungsoperator,
    \item $\Chi$ – stabilisierter Verzweigungsoperator.
\end{itemize}

Die Operatoralgebra wird bis zu den Äquivalenzklassen abelsch:
\[
[\Psi,\Phi]=[\Phi,\Lambda]=\dots = 0.
\]

% ================================================================
\subsubsection{Prozesse auf $K_{12}$}

\begin{itemize}
    \item Konvergenz rekursiver Strukturen,
    \item Abschluss verzweigter Topologien,
    \item Bildung globaler Fixpunkte,
    \item Zusammenbruch unzulässiger Hierarchien,
    \item Entstehung vollständiger Symmetrie unter Operatoren,
    \item Stabilisierung der Zeit- und Zyklusstruktur.
\end{itemize}

% ================================================================
\subsubsection{Vorhersagen für $K_{12}$}

\begin{enumerate}
    \item Jede zulässige rekursive Metahierarchie konvergiert schließlich zu
          Strukturen darstellbar in $K_{12}$.
    \item Über diese Ebene hinaus können keine neuen Achsen oder Potenziale entstehen.
    \item Alle Operatoraktionen reduzieren sich auf globale Symmetrien.
    \item Verzweigung wird topologisch geschlossen; Es entstehen keine neuen Zweige.
    \item Wenn kein globaler Fixpunkt erreicht wird, erzwingt der Kollaps einen Rückfall auf $K_{11}$.
\end{enumerate}

% ================================================================
\subsubsection{Experimente für $K_{12}$}

Es existieren nur indirekte Experimente:
\begin{itemize}
    \item Simulationen rekursiver Konvergenz,
    \item Festkommaberechnung in Metaräumen höherer Ordnung,
    \item Invarianztests unter Operatoralgebra,
    \item Stabilitätstests für Verzweigungsverschlüsse,
    Zulässigkeit der \item-Modellierung unter $M_{12}$.
\end{itemize}

% ================================================================
\subsubsection{Zusammenbruch und Tod von $K_{12}$}

Tod, wenn:
\[
\Omega(K_{12})=\varnothing.
\]

Mechanismen:
\begin{itemize}
    \item rekursive Divergenz auf tiefen Ebenen,
    Verstöße gegen die \item-Operatorinvarianz,
    \item Fehler beim Verzweigungsabschluss,
    \item globale Inkohärenz von Fixpunkten,
    \item Inkompatibilität mit dem Metaraum $M_{12}$.
\end{itemize}

% ================================================================
\subsubsection{Falschbarkeit von $K_{12}$}

Verfälscht, wenn:
\begin{itemize}
    \item eine Metahierarchie konvergiert, ohne invariante Fixpunkte zu bilden,
    Die \item-Verzweigung bleibt nicht geschlossen, die Stabilität bleibt jedoch erhalten.
    \item neue Achsen oder Potenziale entstehen jenseits der Schließung,
    \item-Operatoralgebra wird nicht symmetrisch,
    Die zulässige Rekursion von \item führt zu nicht konvergenten Hierarchien.
\end{itemize}

% ================================================================
\subsubsection{Verzweigung / ontologische Position von $K_{12}$}

Verzweigungsstruktur:
\begin{itemize}
    \item alle Zweige schließen,
    \item Es bilden sich keine neuen Zweige,
    \item die Topologie wird global kompakt und invariant.
\end{itemize}

Ontologische Position:
\[
K_{11} \;\longrightarrow\; K_{12} \quad \text{(Grenzkontinuum)}.
\]

Über $K_{12}$ hinaus postuliert die Theorie keine neuen Ebenen; das Grenzkontinuum
enthält alle zulässigen Strukturmöglichkeiten.

% ================================================================
\subsubsection{Beziehung zu M-Räumen}

$M_{12}$ definiert:
\begin{itemize}
    \item vollständiger Zulässigkeitsbereich für Grenzstrukturen,
    \item-Invarianzanforderungen für Operatoren,
    \item globale Kohärenzbeschränkungen,
    \item Verzweigungsschlussgrenzen,
    \item Fixkomma-Existenzbedingungen,
    \item Kompatibilität mit allen $M_0\dots M_{11}$.
\end{itemize}

Ein Kontinuum existiert auf der Ebene $K_{12}$ genau dann, wenn seine gesamte Struktur innerhalb der Ebene liegt
zulässiger Bereich von $M_{12}$.
